\Needspace{6\baselineskip}
\begin{center}\hypertarget{native-260-article187}{}\pdfbookmark[2]{187.}{article-187}187.\end{center}
Revertimur ab his disquisitionibus generalibus ad propositum nostrum proprium,
cuius caussa illae susceptae fuerant. Antequam determinationem quam
exactissimam orbitae ex observationibus pluribus, quam quot necessario requiruntur,
aggredi liceat, determinatio approximata iam adesse debet, quae ab
omnibus observationibus datis haud multum discrepet. Correctiones his elementis
approximatis adhuc applicandae, ut consensus quam accuratissimus
efficiatur, tamquam problematis quaesita considerabuntur. Quas quum tam
\SourceLeaf{261}exiguas evasuras esse supponi possit, ut quadrata productaque negligere liceat,
variationes, quas corporis coelestis loca geocentrica computata inde nanciscuntur,
per formulas differentiales in Sect. secunda Libri primi traditas computari poterunt.
Loca igitur secundum elementa correcta quae quaerimus computata, exhibebuntur
per functiones lineares correctionum elementorum, illorumque comparatio
cum locis observatis secundum principia supra exposita ad determinationem
valorum maxime probabilium perducet. Hae operationes tanta simplicitate
gaudent, ut ulteriori illustratione opus non habeant, sponteque patet, observationes
quotcunque et quantumvis ab invicem remotas in usum vocari posse. --- Eadem
methodo etiam ad correctionem orbitarum \emph{parabolicarum} cometarum uti licet,
si forte observationum series longior adest, consensusque quam optimus postulatur.
